% GENERATED FILE. Edit canonical lessons, then run npm run generate:books.
% canonical-source-hash: fnv1a64:415cf09e7c11d0ca
% canonical-lessons: KA-C188-gudu, KA-C188-balli, KA-C188-moggu, KA-C188-mullu, KA-C188-tulasi

\chapter{Nest, Creeper, Bud, Thorn, Holy basil}
\label{ch:ka-nest-creeper-bud-thorn-holy-basil}
\hlchaptermodality{188}

\begin{chapteropening}
\textbf{By the end of this chapter:} \emph{I can say a nest, a creeper, a bud, a thorn and holy basil.}

Complete the last lesson of chapter 188: I can say a nest, a creeper, a bud, a thorn and holy basil.
\end{chapteropening}

\section[\texorpdfstring{\kn{ಗೂಡು} (gūḍu)}{ಗೂಡು (gūḍu)}]{\kn{ಗೂಡು} (gūḍu) — a nest}

\label{lesson:KA-C188-gudu}



\begin{quote}
Before the new one: say the Kannada for shade, then the Kannada for darkness.
\end{quote}

\subsection*{You'll want to know: \kn{ಗೂಡು}}
\textbf{\kn{ಗೂಡು}} — \emph{gūḍu} — \textquotedblleft{}a nest\textquotedblright{}.

\begin{grammarlens}[title={What you've built}]
One more word for this chapter.
\end{grammarlens}

\subsection*{Your turn}
\begin{itemize}
\raggedright
  \item \emph{Say it:} \emph{gūḍu}
  \item \emph{Say it:} \emph{gūḍu}, once more
  \item \emph{Say it:} say \emph{gūḍu}
  \item \emph{Recall:} say the Kannada for shade, then the Kannada for darkness, then say \emph{gūḍu} again
\end{itemize}

\begin{tcolorbox}[breakable,colback=teal!4,colframe=teal!35!black,title={Before you move on}]
What does \kn{ಗೂಡು} mean? (\textquotedblleft{}A nest\textquotedblright{}.) Say it once more.
\end{tcolorbox}
\section[\texorpdfstring{\kn{ಬಳ್ಳಿ} (baḷḷi)}{ಬಳ್ಳಿ (baḷḷi)}]{\kn{ಬಳ್ಳಿ} (baḷḷi) — a creeper, a vine}

\label{lesson:KA-C188-balli}



\begin{quote}
Before the new one: say the Kannada for darkness, then the Kannada for a nest.
\end{quote}

\subsection*{You'll want to know: \kn{ಬಳ್ಳಿ}}
\textbf{\kn{ಬಳ್ಳಿ}} — \emph{baḷḷi} — \textquotedblleft{}a creeper, a vine\textquotedblright{}.

\begin{grammarlens}[title={What you've built}]
One more word for this chapter.
\end{grammarlens}

\subsection*{Your turn}
\begin{itemize}
\raggedright
  \item \emph{Say it:} \emph{baḷḷi}
  \item \emph{Say it:} \emph{baḷḷi}, once more
  \item \emph{Say it:} say \emph{baḷḷi}
  \item \emph{Recall:} say the Kannada for darkness, then the Kannada for a nest, then say \emph{baḷḷi} again
\end{itemize}

\begin{tcolorbox}[breakable,colback=teal!4,colframe=teal!35!black,title={Before you move on}]
What does \kn{ಬಳ್ಳಿ} mean? (\textquotedblleft{}A creeper, a vine\textquotedblright{}.) Say it once more.
\end{tcolorbox}
\section[\texorpdfstring{\kn{ಮೊಗ್ಗು} (moggu)}{ಮೊಗ್ಗು (moggu)}]{\kn{ಮೊಗ್ಗು} (moggu) — a bud}

\label{lesson:KA-C188-moggu}



\begin{quote}
Before the new one: say the Kannada for a nest, then the Kannada for a creeper.
\end{quote}

\subsection*{You'll want to know: \kn{ಮೊಗ್ಗು}}
\textbf{\kn{ಮೊಗ್ಗು}} — \emph{moggu} — \textquotedblleft{}a bud\textquotedblright{}.

\begin{grammarlens}[title={What you've built}]
One more word for this chapter.
\end{grammarlens}

\subsection*{Your turn}
\begin{itemize}
\raggedright
  \item \emph{Say it:} \emph{moggu}
  \item \emph{Say it:} \emph{moggu}, once more
  \item \emph{Say it:} say \emph{moggu}
  \item \emph{Recall:} say the Kannada for a nest, then the Kannada for a creeper, then say \emph{moggu} again
\end{itemize}

\begin{tcolorbox}[breakable,colback=teal!4,colframe=teal!35!black,title={Before you move on}]
What does \kn{ಮೊಗ್ಗು} mean? (\textquotedblleft{}A bud\textquotedblright{}.) Say it once more.
\end{tcolorbox}
\section[\texorpdfstring{\kn{ಮುಳ್ಳು} (muḷḷu)}{ಮುಳ್ಳು (muḷḷu)}]{\kn{ಮುಳ್ಳು} (muḷḷu) — a thorn}

\label{lesson:KA-C188-mullu}



\begin{quote}
Before the new one: say the Kannada for a creeper, then the Kannada for a bud.
\end{quote}

\subsection*{You'll want to know: \kn{ಮುಳ್ಳು}}
\textbf{\kn{ಮುಳ್ಳು}} — \emph{muḷḷu} — \textquotedblleft{}a thorn\textquotedblright{}.

\begin{grammarlens}[title={What you've built}]
One more word for this chapter.
\end{grammarlens}

\subsection*{Your turn}
\begin{itemize}
\raggedright
  \item \emph{Say it:} \emph{muḷḷu}
  \item \emph{Say it:} \emph{muḷḷu}, once more
  \item \emph{Say it:} say \emph{muḷḷu}
  \item \emph{Recall:} say the Kannada for a creeper, then the Kannada for a bud, then say \emph{muḷḷu} again
\end{itemize}

\begin{tcolorbox}[breakable,colback=teal!4,colframe=teal!35!black,title={Before you move on}]
What does \kn{ಮುಳ್ಳು} mean? (\textquotedblleft{}A thorn\textquotedblright{}.) Say it once more.
\end{tcolorbox}
\section[\texorpdfstring{\kn{ತುಳಸಿ} (tuḷasi)}{ತುಳಸಿ (tuḷasi)}]{\kn{ತುಳಸಿ} (tuḷasi) — holy basil}

\label{lesson:KA-C188-tulasi}



\begin{quote}
Before the new one: say the Kannada for a bud, then the Kannada for a thorn.
\end{quote}

\subsection*{You'll want to know: \kn{ತುಳಸಿ}}
\textbf{\kn{ತುಳಸಿ}} — \emph{tuḷasi} — \textquotedblleft{}holy basil\textquotedblright{}.

\begin{grammarlens}[title={What you've built}]
One more word for this chapter.
\end{grammarlens}

\subsection*{Your turn}
\begin{itemize}
\raggedright
  \item \emph{Say it:} \emph{tuḷasi}
  \item \emph{Say it:} \emph{tuḷasi}, once more
  \item \emph{Say it:} say \emph{tuḷasi}
  \item \emph{Recall:} say the Kannada for a bud, then the Kannada for a thorn, then say \emph{tuḷasi} again
\end{itemize}

\begin{tcolorbox}[breakable,colback=teal!4,colframe=teal!35!black,title={Before you move on}]
What does \kn{ತುಳಸಿ} mean? (\textquotedblleft{}Holy basil\textquotedblright{}.) Say it once more.
\end{tcolorbox}
